\documentclass[10pt,a4paper]{amsart}
\usepackage[margin=19mm]{geometry}
\usepackage{amsmath,amssymb,mathtools}
\usepackage[scr=rsfs,cal=euler]{mathalfa}
\usepackage{xfrac}
\usepackage[unicode,hidelinks,bookmarks=false]{hyperref}
\newcommand{\ie}{i.e.\ }
\newcommand{\cf}{cf.\ }
\newcommand{\resp}{respectively\ }
\newcommand{\Subsection}{\subsection}
\providecommand{\nobreakdash}{\nobreak\hbox{-}}
\setlength{\emergencystretch}{2em}
\multlinegap0pt
\usepackage{bm}
\usepackage{bbm}
\usepackage{dsfont}
\usepackage{xspace}
\usepackage{cancel}
\usepackage{mathtools}
\usepackage{cleveref}
\usepackage{amsthm}
\makeatletter
\@ifundefined{theorem}{\newtheorem{theorem}{Theorem}}{}
\@ifundefined{lemma}{\newtheorem{lemma}[theorem]{Lemma}}{}
\@ifundefined{proposition}{\newtheorem{proposition}[theorem]{Proposition}}{}
\makeatother
\makeatletter
\newcommand{\Argmax}{\mathop{\rm argmax}}
\newcommand{\T}{\mathcal{T}}
\newcommand{\cA}{\mathcal{A}}
\newcommand{\cS}{\mathcal{S}}
\newcommand*{\infn}[1]{\left\|{#1}\right\|_{\infty}}
\expandafter\ifx\csname manualcorollary\endcsname\relax
\newenvironment{manualcorollary}[1]{%
    \renewcommand\themanualcorollaryinner{#1}%
  \manualcorollaryinner
}{\endmanualtheoreminner}
\fi
\expandafter\ifx\csname manualproposition\endcsname\relax
\newenvironment{manualproposition}[1]{%
  \renewcommand\themanualpropinner{#1}%
  \manualpropinner
}{\endmanualpropinner}
\fi
\expandafter\ifx\csname manualtheorem\endcsname\relax
\newenvironment{manualtheorem}[1]{%
    \renewcommand\themanualtheoreminner{#1}%
  \manualtheoreminner
}{\endmanualtheoreminner}
\fi
\newcommand{\reals}{\mathbb{R}}
\renewcommand{\sp}[1]{\|{#1}\|_\text{\rm sp}}
\makeatother
\title[Near-Optimal Sample Complexity for MDPs via Anchoring]{Near-Optimal Sample Complexity for MDPs via Anchoring}
\author{Jongmin Lee, Mario Bravo, Roberto Cominetti}
\date{Source: arXiv:2502.04477v2 (CC BY 4.0)}
\begin{document}
\maketitle

\noindent\emph{Source and licence.} Near-Optimal Sample Complexity for MDPs via Anchoring. Version arXiv:2502.04477v2 (CC BY 4.0), distributed under the Creative Commons Attribution 4.0 International licence, as recorded in the arXiv licence metadata for this exact version and shown on the abstract page https://arxiv.org/abs/2502.04477v2. Full rights notice: licenses/arxiv-2502.04477-CC-BY-4.0.txt.\par\smallskip

\noindent\textit{Editorial scope.} Counted unit: the Theorem whose source label is \texttt{thm:err-d} in the article (source0.tex lines 1306--1315, source label \texttt{thm:err-d}), delivered together with the article's own complete original proof of it (source0.tex lines 1316--1346, one \texttt{proof} environment of 31 lines). No mathematics is written, repaired or re-derived: statement and proof are the article's own text line for line. Required context retained uncounted: thm:err (lines 531--540); prop:err-d (lines 709--711); prop:sam-err0-d (lines 1247--1251); Eq:Ek-d (lines 1259--1261); Eq:spk-d (lines 1266--1267); Le:uno-d (lines 1269--1271); Le:dos-d (lines 1284--1286). Every definition, display and label of those lines is retained unchanged. Omitted from this package and not redistributed: 1--530, 541--708, 712--1246, 1252--1258, 1262--1265, 1268--1268, 1272--1283, 1287--1305, 1347--1539. Nothing inside a retained excerpt is omitted or altered: every retained body is delivered exactly as the article's lines are written, so no omission receipt of any kind accompanies this package. Citations: the retained excerpts carry no citation command at all, so no bibliography entry of the article is transcribed and no reference is redistributed. Reference adaptation: none was needed and none was made; every reference inside the delivered unit resolves inside it, so no unresolved reference or citation is produced. Printed slips kept literal: no printed word, symbol, hypothesis or number of the counted statement or of its proof was altered. Layout and class: the article's own class is \texttt{article} with packages including microtype, graphicx, subfigure, booktabs, amssymb, pifont, hyperref, verbatim, icml2025, amsmath, mathtools, amsthm, cleveref, todonotes; the wrapper is the lane's pinned \texttt{amsart} editorial wrapper at 10pt on A4, which restates the theorem-like environments the retained text needs and adds the article's own macro definitions that the retained text needs (\texttt{\textbackslash Argmax}, \texttt{\textbackslash T}, \texttt{\textbackslash cA}, \texttt{\textbackslash cS}, \texttt{\textbackslash infn}, \texttt{\textbackslash reals}, \texttt{\textbackslash sp}), with extra wrapper packages (bm, bbm, dsfont, xspace, cancel, mathtools, cleveref). The delivered numbering is the unit's own. Assets: no retained span contains a figure environment, a graphics-inclusion command, a PDF-page inclusion, a table or a TikZ picture; no image file is copied into this package, the package declares no asset dependency and no image file is redistributed with it. Licence: version arXiv:2502.04477v2 of this article is distributed under the Creative Commons Attribution 4.0 International licence, as recorded in the arXiv licence metadata for this exact version and shown on the abstract page https://arxiv.org/abs/2502.04477v2. A full rights notice is delivered at licenses/arxiv-2502.04477-CC-BY-4.0.txt. Characters: every retained excerpt is pure ASCII, so the delivered engine, XeLaTeX with its default Latin Modern text font, reports no missing character.
\begin{theorem}\label{thm:err}
Assume \mbox{\sc (H)} and \mbox{\sc (S)} and let $(Q^n,T^n,\pi^n)$ be the output computed by $\mbox{\rm SAVIA}(Q^0,n,\varepsilon,\delta)$. Then, with probability at least $(1-\delta)$ we have, for all $s \in \cS$
\vspace{-1ex}
$$ g^*-g_{\pi^n}(s)\! \leq \sp{Q^n-\T(Q^n)}\leq \frac{8\sp{Q^0-Q^*}}{n+2}+4\varepsilon,$$

\vspace{-2ex}
with a sample and time complexity of order $${O}\big(L\,|\cS||\cA|\big(({\sp{Q^0-Q^*}^{2}+\sp{Q^0}^{2}})/{\varepsilon^{2}}+n^2\big)\big)$$ 
where  
$L=\ln(|\cS||\cA|(n+1)/\delta)\ln^3(n+2)$.
\end{theorem}

\begin{proposition} \label{prop:err-d}
    Let $Q\in\reals^{\cS\times\cA}$ and $\pi:\cS\to\cA$ a greedy policy such that $\pi(s)\in\Argmax_{a\in\cA}~Q(s,a)$. Then we have $ \infn{Q^*-Q_{\pi}}\leq \frac{2}{1-\gamma}\infn{Q-\T_\gamma(Q)}$.
\end{proposition}

\begin{proposition}\label{prop:sam-err0-d}
Let $c_k>0$  with \mbox{$2\sum_{k=0}^\infty c_k^{-1}\leq 1$} and $T^k, Q^k$ the iterates generated by \mbox{\rm SAVID}$(Q^0,n,\varepsilon,\delta, \gamma)$.
Then, with probability at least $1-\delta$ we have $\|T^k-\T_\gamma(Q^k)\|_\infty\le  \gamma \varepsilon$ simultaneously for 
 all $k=0,\ldots,n$.
\end{proposition}

\begin{equation}\label{Eq:Ek-d}
\infn{T^k-\T_\gamma(Q^{k})}\le \varepsilon\qquad \text{for all} \,\, k=0,\ldots,n.
\end{equation}

\begin{equation}\label{Eq:spk-d}\infn{T^k-T^{k-1}}=\gamma \infn{D^k} \le \gamma \infn{d^k}=\gamma \infn{V^k-V^{k-1}}\le \gamma \infn{Q^{k}-Q^{k-1}}.
\end{equation}

\begin{lemma}\label{Le:uno-d}Assuming \mbox{\sc (S)} and \eqref{Eq:Ek-d}, we have that 
$$\infn{Q^k-Q^*}\leq \infn{Q^0-Q^*}+\frac{1}{3}\,\varepsilon\,k \quad \text{for all} \,\, k=0,\ldots,n.$$
\end{lemma}

\begin{lemma}\label{Le:dos-d} Assume \mbox{\sc (S)} and \eqref{Eq:Ek-d} and let $\rho_k=2\infn{Q^0-Q^*}+\frac{1}{3}\,\varepsilon\,k$. Then 
$$\infn{Q^{k}-Q^{k-1}} \le \frac{2}{k(k+1)}\sum_{i=1}^{k}\rho_{i+2}.\quad \quad \text{for all} \,\, k=1,\ldots,n.$$
\end{lemma}

\begin{theorem}\label{thm:err-d}
Assume \mbox{\sc (S)} and let $(Q^n,T^n,\pi^n)$ be the output computed by $\mbox{\rm SAVID}(Q^0,n,\varepsilon,\delta, \gamma)$. Then, with probability at least $(1-\delta)$ we have
\vspace{-1ex}
$$  \infn{Q^n-\T_\gamma(Q^n)}\leq \frac{8\infn{Q^0-Q^*}}{n+2}+2\varepsilon,$$

\vspace{-2ex}
with a sample and time complexity of order $${O}\big(L_\gamma\,|\cS||\cA|\big(({\infn{Q^0-Q^*}^{2}+\infn{Q^0}^{2}})/{\varepsilon^{2}}+n^2\big)\big)$$ 
where  
$L_\gamma=\ln(2|\cS||\cA|(n+1)/\delta)\ln^3(n+2)$.
\end{theorem}

\begin{proof}
From \cref{prop:sam-err0-d}, with probability at least $(1-\delta)$ we have $\infn{T^k-\T_\gamma(Q^k)}\leq\varepsilon$ for all $k=0,\ldots,n$. Now, the recursion $Q^n=\frac{2}{n+2}Q^0+\frac{n}{n+2}T^{n-1}$ implies
 \begin{equation}\label{eq:lll-d}
    Q^n-\T_\gamma(Q^n)
    =\mbox{$\frac{2}{n+2}\big(Q^0-\T_\gamma(Q^n)\big)+ \frac{n}{n+2}\big(T^{n-1}-\T_\gamma(Q^{n-1})\big)+ \frac{n}{n+2} \big(\T_\gamma(Q^{n-1})-\T_\gamma(Q^n)\big)$}.
\end{equation}
Using the triangle inequality and Lemma~\ref{Le:uno-d}, we have that  
$$\infn{Q^0-\T_\gamma(Q^n)} \le \infn{Q^0-Q^*}+\gamma \infn{Q^*-Q^n} \le 2\infn{Q^0-Q^*}+\mbox{$\frac{1}{3}$}\,\varepsilon\,n =  \rho_n,$$ 
while \cref{Le:dos-d} gives $\infn{\T_\gamma(Q^{n-1})-\T_\gamma(Q^n)}\leq\frac{2}{n(n+1)}\sum_{i=1}^{n}\rho_{i+2}$. Thus, applying a triangle inequality to \eqref{eq:lll-d} and using these estimates together with \cref{prop:err-d}, we obtain 
\begin{align*}
    \infn{Q^n-\T_\gamma(Q^n)}&\le \mbox{$\frac{2}{n+2}\,\rho_n+ \varepsilon+\frac{2}{(n+1)(n+2)}\sum_{i=1}^{n}\rho_{i+2}$}
   \\& = \mbox{$\frac{4(1+2n)}{(n+1)(n+2)}\infn{Q^0-Q^*}+ \frac{2(3+8n+3n^2)}{3(n+1)(n+2)}\varepsilon$}
    \\& \le \mbox{$\frac{8\infn{Q^0-Q^*}}{n+2}+2\varepsilon$}.
\end{align*}
To estimate the complexity, for $k \ge 1$, we use the inequality $\infn{d^k}\le\infn{Q^k-Q^{k-1}}$ in \eqref{Eq:spk-d} together with \cref{Le:dos-d} to find
\begin{align*}
    \mbox{$\infn{d^k}$}&\leq\mbox{$\frac{2}{k(k+1)}\sum_{i=1}^{k}\rho_{i+2}$}\\
    &=\mbox{$\frac{4}{k+1}\infn{Q^0-Q^*}$} +\mbox{$\frac{k+5}{3(k+1)}\varepsilon$}\\
    &\leq \mbox{$\frac{4}{k+1}\infn{Q^0-Q^*}+\varepsilon$}.
\end{align*}
Now, to estimate the total number of samples $|\cS||\cA|\sum_{k=0}^n m_k$ 
we recall that $m_k=\max\{\lceil 2 \alpha c_k\infn{d^{k}}^2/\varepsilon^{2}\rceil,1\}$ which can be bounded as $m_k\leq 1+2\alpha c_k\infn{d^{k}}^2/\varepsilon^{2}$. Then 
\begin{align}
    \mbox{$\sum_{k=0}^nm_k$}&\leq\mbox{$ (n+1)+\frac{2\alpha}{\varepsilon^{2}}\sum_{k=0}^n c_k\infn{d^{k}}^2$} \nonumber\\
   &\leq\mbox{$ (n+1)+\frac{20\, \alpha}{\varepsilon^{2}} \ln^2(2)\,\infn{Q_0}^2+\frac{10\alpha}{\varepsilon^{2}}\sum_{k=1}^n (k+2)\ln^2(k+2)\big(\frac{4}{k+1}\infn{Q^0-Q^*}+\varepsilon\big)^2$} \nonumber\\
      &\leq\mbox{$ (n+1)+\frac{20\, \alpha}{\varepsilon^{2}} \ln^2(2)\infn{Q_0}^2+\sum_{k=1}^n\frac{480\,\alpha}{\varepsilon^{2}(k+2)}\ln^2(k+2)\infn{Q^0-Q^*}^2+20\,\alpha\sum_{k=1}^n(k+2)\ln^2(k+2)$} \nonumber\\
   &=O\!\left( \alpha \infn{Q^0}^2/\varepsilon^{2}+\mbox{$\alpha\ln^3(n+2)\infn{Q^0-Q^*}^2/\varepsilon^{2}+\alpha \,n^2\ln^2(n+2)$}\right),
   \label{eq:last-d}
\end{align}
where we use the same estimations as in the proof of \cref{thm:err}.
\end{proof}

\end{document}
